% Regression: t2_sptmpo_cross — pulls an MPO symmetry string through an SPT PEPS tensor and deposits U_g.
% Formula: (schematic, per site):
% T2 — RMP arXiv:2011.12127, fig. fig3-pepe_newSPTMPO: pulling the
% symmetry MPO string through an SPT PEPS tensor.
% Formula (schematic, per site):
%   [g-string crossing the W and NE virtual legs ABOVE the tensor] . A
%   = (U_g on the physical leg) . A . [g-string crossing the SW and E
%      virtual legs BELOW the tensor],
% i.e. sliding the red MPO string g across the PEPS tensor deposits the
% on-site action U_g on the physical index; at each crossed virtual leg
% the string carries a crossing tensor (source ink: small parallelogram).
% Free tier (genuinely non-grid: a four-valent star with diagonal legs
% and curved string arcs).
%
% Missing primitives (checked in tenkz-free.code.tex):
% 1. \tnjoin strokes `bond`/`fused bond` only -- no operator-hued join
%    (no /tenkz/join/operator key), so the g-string the source draws RED
%    renders black.  Load-bearing: hue is what separates the MPO string
%    from the black lattice legs.
% 2. No sheared/parallelogram crossing-tensor skin in \tnput
%    (dot|box|pill|mpo|ring only); upright boxes stand in.
\documentclass[varwidth=19cm, border=6pt]{standalone}
\usepackage{amsmath, amssymb}
\usepackage{tenkz}
\begin{document}
\[
  % LHS: the string crosses W and NE legs, passing ABOVE the tensor
  \begin{tenkzfree}
    \tnput[dot]{c}{(0,0)}{}
    \tnput[box]{p1}{(-0.75,0)}{}
    \tnput[box]{p2}{(0.62,0.52)}{}
    % lattice legs: W--E in-plane pair and the SW--NE diagonal pair
    \tnjoin{-1.5,0}{p1.west}
    \tnjoin{p1.east}{c}
    \tnjoin{c}{1.5,0}
    \tnjoin{-1.1,-0.92}{c}
    \tnjoin{c}{p2.south west}
    \tnjoin{p2.north east}{1.25,1.05}
    % physical leg
    \tnjoin{c}{0,0.6}
    % the g-string (source: RED): in from the lower west, over the tensor
    \tnjoin[route=curve, out=35, in=200]{-1.55,-0.5}{p1}
    \tnjoin[route=curve, out=55, in=195, label=g]{p1}{p2}
    \tnjoin[route=curve, out=15, in=185]{p2}{1.65,0.8}
  \end{tenkzfree}
  \;=\;
  % RHS: U_g on the physical leg; the string passes BELOW the tensor
  \begin{tenkzfree}
    \tnput[dot]{c}{(0,0)}{}
    \tnput[dot, label pos=west]{u}{(0,0.55)}{U_g}
    \tnput[box]{p3}{(0.75,0)}{}
    \tnput[box]{p4}{(-0.62,-0.52)}{}
    % lattice legs
    \tnjoin{-1.5,0}{c}
    \tnjoin{c}{p3.west}
    \tnjoin{p3.east}{1.5,0}
    \tnjoin{-1.25,-1.05}{p4.south west}
    \tnjoin{p4.north east}{c}
    \tnjoin{c}{1.1,0.92}
    % physical leg with the deposited on-site action U_g
    \tnjoin{c}{u}
    \tnjoin{u}{0,1.1}
    % the g-string (source: RED): in from the west, under the tensor
    \tnjoin[route=curve, out=0, in=145]{-1.55,-0.35}{p4}
    \tnjoin[route=curve, out=-35, in=245, label=g]{p4}{p3}
    \tnjoin[route=curve, out=35, in=200]{p3}{1.65,0.65}
  \end{tenkzfree}
\]
\end{document}
